\chapter{User-Side Job Recommendation System}
\label{chap:user_recommendation}

This chapter describes the job recommendation system built for users, consisting of the data pipeline and the two-tower neural network architecture. The system learns dense representations of both users and jobs and uses their similarity to rank job opportunities for each individual user. This model is later used in Chapter~\ref{chap:congestion_balancing} as the basis for the user-side congestion reduction strategy.

A two-tower \gls{RS} is chosen because neural networks can handle the cold-start problem~\cite{rama2019deeplearningaddresscandidate} — a user with no prior interaction history can still receive recommendations based on their profile features alone by using the learned user embedding. This property is essential given the temporal evaluation setup described in Chapter~\ref{chap:experiments}.

Section~\ref{sec:user_arch} describes the architecture of the model. Section~\ref{sec:user_data_pipeline} describes the data pipeline and feature engineering. Section~\ref{sec:user_training} describes the training procedure. All experimental results and hyperparameter choices are reported in Chapter~\ref{chap:experiments}.

\section{Architecture: The Two-Tower Model}
\label{sec:user_arch}

The two-tower model consists of two separate neural networks. The user tower processes user features and the job tower processes job features. Both towers learn a dense vector representation of its respective input. The compatibility of a user-job pair is then measured as the dot product of these two vectors, as defined in Section~\ref{sec:notation}. The design allows job embeddings to be precomputed and cached, enabling efficient real-time recommendations for users. 

\subsection{User Tower}
\label{sec:user_tower}

The user tower is a multi-layer feedforward network that transforms the input user feature vector $\mathbf{x}_u$ into a $d$-dimensional embedding $\mathbf{e}_u$. Categorical input features are first passed through learnable embedding layers (see Section~\ref{sec:prelim_embeddings}), producing dense representations for each categorical attribute. These are concatenated with the numerical features to form the full input vector. This vector is then processed true a sequence of fully-connected layers with batch normalisation, ReLU activations and dropout regularisation.

The layer progression narrows from a wide hidden representation down to the final $d$-dimensional output, acting as a bottleneck that forces the network to learn a compressed and meaningful summary of the user profile. The architecture is detailed in Algorithm~\ref{alg:user_tower}.

\begin{algorithm}[H]
\caption{User Tower Forward Pass}
\label{alg:user_tower}
\begin{algorithmic}[1]
\Require Categorical features $\mathbf{c}_u$, numerical features $\mathbf{n}_u$
\For{each categorical feature $c_k$ in $\mathbf{c}_u$}
    \State $\mathbf{emb}_k \gets \text{Embedding}_{k}(c_k)$
\EndFor
\State $\mathbf{x}_u \gets [\mathbf{emb}_1, \ldots, \mathbf{emb}_K, \mathbf{n}_u]$ 
\State $\mathbf{h}_1 \gets \text{Dropout}(\text{ReLU}(\text{BN}(\mathbf{W}_1 \mathbf{x}_u + \mathbf{b}_1)))$
\State $\mathbf{h}_2 \gets \text{Dropout}(\text{ReLU}(\text{BN}(\mathbf{W}_2 \mathbf{h}_1 + \mathbf{b}_2)))$
\State $\mathbf{e}_u \gets \mathbf{W}_3 \mathbf{h}_2 + \mathbf{b}_3$
\Ensure User embedding $\mathbf{e}_u \in \mathbb{R}^d$
\end{algorithmic}
\end{algorithm}

\subsection{Job Tower}
\label{sec:job_tower}

The job tower mirrors the user tower architecture exactly, processing job features $\mathbf{x}_j$ to produce a $d$-dimensional job embedding $\mathbf{e}_j$. The same pattern of categorical embeddings concatenated with numerical features, followed by the bottleneck \gls{MLP}, is applied. The only structural difference is the input dimension, which depends on the number of job-specific features. This symmetry is intentional: it ensures that user and job embeddings live in the same $d$-dimensional space, making the dot product a meaningful similarity measure.

\subsection{Prediction Layer}
\label{sec:user_prediction}

The final user-job compatibility score is computed as the dot product of the two learned embeddings:
\begin{equation}
s(\mathbf{e}_u, \mathbf{e}_j) = \mathbf{e}_u \cdot \mathbf{e}_j = \sum_{i=1}^{d} (\mathbf{e}_u)_i \cdot (\mathbf{e}_j)_i
\end{equation}

Higher scores indicate stronger predicted affinity between a user's preferences and a job's characteristics. At training time, this raw score is passed through a sigmoid to produce a probability for the binary cross-entropy loss because the label is binary. At inference time, the raw score is used directly for ranking, as discussed in Section~\ref{sec:inference_postprocessing}.

L2 normalisation of the output embeddings was considered as an alternative design. The impact of this choice on model performance is evaluated in the ablation study in Section~\ref{sec:exp_user_ablation}.

\section{Data Pipeline and Feature Engineering}
\label{sec:user_data_pipeline}

The raw data must be preprocessed before it can be used by the two-tower model.The pipeline consists of two main stages: user feature processing and job feature processing. Both stages include text embedding as part of their respective pipelines. Each stage is designed to handle large-scale data efficiently. The construction of the training dataset, including the negative sampling strategy, is described in Chapter~\ref{chap:experiments}.



\subsection{User Feature Processing}
\label{sec:user_features_processing}

User features are extracted and engineered to form a comprehensive representation of each job seeker's profile.

\subsubsection{Feature Extraction and Merging}

Raw user features — demographics, education and location — are merged with the user's employment history. The most recent job titles (up to three previous positions) are extracted and concatenated with recency labels to form a single career trajectory string:
\begin{equation}
\text{Recent\_JobTitle} = \text{``Current: } t_1 \text{ | Previous: } t_2 \text{ | Earlier: } t_3\text{''}
\end{equation}
where $t_1, t_2, t_3$ are the three most recent job titles in reverse chronological order. This string is subsequently embedded using the Qwen3 text embedding model to capture the semantic meaning of the user's career progression.

\subsubsection{Feature Engineering}

Several domain-specific features are constructed from raw attributes. Years of professional experience is derived from the graduation date as $\text{current\_year} - \text{graduation\_year}$, since elapsed time since graduation is a more informative career-stage signal than the raw date. Missing values are imputed with the median. Location fields are standardised by retaining city and state while dropping redundant or noisy columns, ensuring a consistent geographic representation across all users.

\subsubsection{Encoding}
\label{sec:user_embedding}

Each feature type is processed according to its nature. Numerical features such as years of experience are standardised using z-score normalisation to ensure scale invariance at the MLP input. Categorical features such as education level are label-encoded to integers and fed into learnable embedding layers rather than the MLP directly, so no further normalisation is required. The career trajectory string is encoded into a $d_{\text{text}}$-dimensional dense vector using the Qwen3 Embedding model, capturing semantic relationships between job titles and career stages. The embedding dimension $d_{\text{text}}$ is chosen so that the text signal contributes meaningfully but does not dominate the other features; this choice is discussed in Chapter~\ref{chap:experiments}.

Table~\ref{tab:user_features} summarises the full user feature engineering pipeline.

\begin{table}[htbp]
    \centering
    \caption{User Feature Engineering Pipeline}
    \label{tab:user_features}
    \begin{tabular}{|l|l|l|l|l|}
        \hline
        \textbf{Category} & \textbf{Raw Feature} & \textbf{Type} & \textbf{Processing} & \textbf{Output} \\
        \hline
        \multirow{2}{*}{\textbf{Identifiers}} & UserID & int & --- & UserID \\
        & WindowID & int & --- & WindowID \\
        \hline
        \multirow{2}{*}{\textbf{Location}} & City & categorical & LabelEncoder & int [1, N] \\
        & State & categorical & LabelEncoder & int [1, N] \\
        \hline
        \multirow{3}{*}{\textbf{Education}} & DegreeType & categorical & LabelEncoder & int [1, N] \\
        & Major & categorical & LabelEncoder & int [1, N] \\
        & GraduationDate & date & $\rightarrow$ YearsSinceGrad & float (scaled) \\
        \hline
        \multirow{5}{*}{\textbf{Work Experience}} & WorkHistoryCount & int & StandardScaler & float (scaled) \\
        & TotalYearsExperience & int & StandardScaler & float (scaled) \\
        & CurrentlyEmployed & categorical & LabelEncoder & int [0, 1] \\
        & ManagedOthers & categorical & LabelEncoder & int [0, 1] \\
        & ManagedHowMany & int & StandardScaler & float (scaled) \\
        \hline
        \textbf{Job History} & Recent\_JobTitle (top 3) & text & Qwen Embedding & $\mathbb{R}^{d_{\text{text}}}$ \\
        \hline
    \end{tabular}
\end{table}

\subsection{Job Feature Processing}
\label{sec:job_features_processing}

Job features are processed using a two-pass chunked strategy to handle the large scale of the job dataset without loading it entirely into memory.

In the first pass, all records are iterated to collect unique categorical values and compute scaling statistics — mean and variance — for numerical features across the full dataset. In the second pass, these fitted encoders and scalers are applied consistently, ensuring that categorical mappings and numerical normalisation are coherent across all jobs. Numerical features are standardised using z-score normalisation; categorical features are mapped to integers using a consistent encoding dictionary (1-indexed, with 0 reserved for missing values).

\subsection{Text Cleaning and Embedding}
\label{sec:text_embedding}

Job titles and descriptions often contain HTML markup, special characters and inconsistent formatting. Text is cleaned by removing HTML tags, lowercasing and stripping whitespace before embedding.

The cleaned job titles are encoded into dense vectors using the Qwen3 Embedding model. As with user job history, the embedding dimension is truncated to $d_{\text{text}}$ to balance the contribution of text against other input features. The embeddings are unit-normalised to ensure comparable similarity scores across the job corpus. These Qwen embeddings capture semantic relationships between job titles --- for example, ``Software Engineer'' and ``Developer'' will have nearby representations --- enabling the job tower to learn meaningful job-to-job relationships beyond categorical attributes alone.

\subsection{Cross-Tower Location Consistency}
\label{sec:location_encoding}

An important preprocessing step is ensuring that the same location — city or state — maps to the same integer index in both the user and job feature files. Without this, the two feature files are encoded independently, meaning a city encoded as integer 42 in the user file may refer to a completely different city than integer 42 in the job file. The shared embedding lookup would then learn no meaningful location signal.

To resolve this, a shared vocabulary is constructed jointly from all city and state values across both datasets, and both files are re-encoded using this common mapping. The practical impact of this step is confirmed in the ablation study (Section~\ref{sec:exp_user_ablation}): including location features without the shared vocabulary hurts performance, while including them with the shared vocabulary yields a substantial improvement in validation AUC. This makes cross-tower location consistency a critical preprocessing requirement for any dataset where location appears in both user and item features.

\section{Training}
\label{sec:user_training}

\subsection{Loss Function}

The model is optimised using Binary Cross-Entropy with Logits loss, which combines a sigmoid activation with binary cross-entropy in a numerically stable form:
\begin{equation}
\mathcal{L} = -\frac{1}{N} \sum_{i=1}^{N} \left[ y_i \log(\sigma(\hat{y}_i)) + (1 - y_i) \log(1 - \sigma(\hat{y}_i)) \right]
\end{equation}
where $y_i \in \{0, 1\}$ is the ground truth label and $\hat{y}_i = s(\mathbf{e}_u, \mathbf{e}_j)$ is the raw dot product score. The sigmoid $\sigma$ maps the unbounded score to a probability during training. The model is optimised using the Adam optimiser with early stopping based on validation loss to prevent overfitting. Specific hyperparameter values are reported in Section~\ref{sec:exp_user_setting}.

\subsection{Inference and Post-Processing}
\label{sec:inference_postprocessing}

At inference time, recommendations are generated efficiently by precomputing all job embeddings once and caching them in memory. For each user, only the user tower forward pass is executed, followed by a single matrix multiplication against the precomputed job matrix. This reduces per-user recommendation time to one matrix multiply regardless of the size of the job pool.

At inference, raw dot product scores are used for ranking rather than sigmoid probabilities. This is because the sigmoid saturates near 1.0 in float32 precision for high-confidence matches, making high-scoring jobs indistinguishable from one another, while the raw scores remain separable.

Geographic preference is applied as a post-processing step rather than being learned by the model. Jobs in the same state as the user receive a ranking boost, ensuring geographically relevant recommendations without conflating location preference with semantic job-fit during training. This separation of concerns — semantic relevance learned by the model, geographic preference applied post-hoc — allows each component to be tuned and evaluated independently. Recommendations are further restricted to jobs from the active evaluation window to avoid surfacing positions that are no longer open.